\begin{table}[ht]
\centering
\caption{Lee bounds on the under-the-gun price gap.}
\label{tab:utg_lee_bounds}
\begin{threeparttable}
\small
\setlength{\tabcolsep}{4pt}
\begin{adjustbox}{max width=\textwidth}
\begin{tabular}{p{.38\linewidth}rrrr}
\toprule
Specification & Admin coef. & SE & Gap (\%) & $N$ \\
\midrule
Naive UTG: item + year + PBU FE & \BPutgPointNaiveCoef{} & \BPutgPointNaiveSE{} & \BPutgPointNaive{} & \BPutgPointNaiveN{} \\
Lee lower bound: admin top-tail trim & \BPutgBoundLowCoef{} & \BPutgBoundLowSE{} & \BPutgBoundLow{} & \BPutgBoundLowN{} \\
Lee upper bound: admin bottom-tail trim & \BPutgBoundHighCoef{} & \BPutgBoundHighSE{} & \BPutgBoundHigh{} & \BPutgBoundHighN{} \\
\bottomrule
\end{tabular}
\end{adjustbox}
\begin{tablenotes}[flushleft]\footnotesize
\item \textit{Notes:} Coefficients are administrative minus litigated log negotiated prices. Negative coefficients mean litigated purchases are more expensive. Percentage gaps are reported as litigated-over-administrative price gaps. Trimming is applied within item $\times$ year $\times$ PBU strata where administrative observations exceed litigated observations. The mean trimming rate is \BPutgLeeTrimMean{}; the maximum trimming rate is \BPutgLeeTrimMax{}.
\end{tablenotes}
\end{threeparttable}
\end{table}
